\section{Riemann Integral}

\begin{definition}[Riemann integral]%
\label{def:integrale}%
\mymark{***}%
\index{integral!definition}%
\index{definition!of integral}%
\index{Riemann!integral of}%
Let $a,b\in \RR$, $a \le b$.
A set $P\subset [a,b]$ is said to be a \emph{Riemann partition}%
\mymargin{Riemann partition}%
\index{partition!of Riemann}
\index{Riemann!partition of}
\index{Riemann partition}
\index{Riemann!partition of}
of the interval $[a,b]$ if $P$ is a finite set such that $a,b\in P$.
In particular, $P$ can be written as
\[
 P = \ENCLOSE{ x_0, x_1, \dots, x_N}
\]
with
\[
  a = x_0 < x_1 < \dots < x_{N-1} < x_N = b.
\]

Let $f\colon [a,b] \to \RR$ be a bounded function.
Given any partition $P$ of $[a,b]$, we define
the \emph{upper sums} and the \emph{lower sums}%
\mymargin{upper/lower sums}%
\index{upper/lower sums}
\index{Riemann!upper sums}
\index{Riemann!lower sums}
as
\begin{align*}
S^*(f,P)
&= \sum_{k=1}^N (x_k - x_{k-1}) \cdot \sup f([x_{k-1},x_k]) \\
S_*(f,P)
&= \sum_{k=1}^N (x_k - x_{k-1}) \cdot \inf f([x_{k-1},x_k]).
\end{align*}
Finally, we define
\begin{align*}
  I^*(f) &= \inf \ENCLOSE{S^*(f,P) \colon \text{$P$ partition of $[a,b]$}}
  \\
  I_*(f) &= \sup \ENCLOSE{S_*(f,P) \colon \text{$P$ partition of $[a,b]$}}.
\end{align*}

If $I^*(f) = I_*(f)$, we say that $f$ is
\emph{Riemann-integrable}%
\mymargin{Riemann integral}%
\index{integral!of Riemann}%
\index{Riemann!integral of}%
\index{integrability}%
\index{integral}%
\index{integral!definition}%
and we say that the \emph{integral} of $f$ over $[a,b]$ is
the common value $I^*(f)=I_*(f)$, which will be denoted by
\[
  \int_a^b f
  \qquad{\text{or by}} \qquad
  \int_a^b f(x)\, dx.
\]

If $b<a$ and if $f$ is Riemann-integrable on $[b,a]$,
we define by convention:
\[
  \int_a^b f = -\int_b^a f.
\]
\end{definition}

\begin{lemma}
Let $f\colon[a,b]\to \RR$ be a bounded function.
If $P$ and $Q$ are any two partitions of the interval
$[a,b]$, then
\begin{equation}\label{eq:5112309}
  S_*(f,P) \le S_*(f,P\cup Q)\le S^*(f,P\cup Q) \le S^*(f,Q).
\end{equation}
Consequently, $I_*(f) \le I^*(f)$.
\end{lemma}
\begin{proof}
    Let $P$ be any partition of $[a,b]$ and let $y\in [a,b]$ be any point.
  Setting $P' = P \cup \ENCLOSE{y}$, we want to show
  that
  \begin{equation}\label{eq:39543}
    S_*(f,P) \le S_*(f,P') \le S^*(f,P') \le S^*(f,P).
  \end{equation}
  If $y\in P$, there is nothing to prove since
    it would result in $P'=P$, and the inequality $S_*(f,P') \le S^*(f,P')$ is
  always satisfied since every supremum appearing in the definition of $S^*$ is
  greater than or equal to the corresponding
  infimum appearing in the definition of $S_*$.
    Suppose then that $y \not \in P$ and thus that $y$ is between two
  consecutive points $x_{k-1}, x_k$ of the partition $P$:
  \[
    a= x_0 < x_1 < \dots < x_{k-1} < y < x_k < \dots < x_N=b.
  \]
  Then the sums defining $S_*(f,P)$ and $S_*(f,P')$ 
  differ only on the interval $[x_{k-1},x_k]$. 
  % and we have
  % \begin{align*}
  %   S_*(f,P') - S_*(f,P)
  %  &= (y-x_{k-1})\cdot \!\!\inf_{[x_{k-1},y]}\!\!\! f
  %  + (x_k - y)\cdot\! \inf_{[y,x_k]}\! f\\
  %  &\quad - (x_k - x_{k-1})\cdot \!\!\inf_{[x_{k-1}, x_k]}\!\!\!f
  % \end{align*}
  Observing that
  \[
  \inf f([x_{k-1}, x_k])
  \le \inf f([x_{k-1},y])
  \qquad \text{and} \qquad
  \inf f([x_{k-1}, x_k])
  \le \inf f([y,x_k])
  \]
  we obtain 
  \begin{align*}
    (x_k - x_{k-1})\cdot \inf f([x_{k-1}, x_k])
    & = (y - x_{k-1} + x_k - y)\cdot \inf f([x_{k-1}, x_k]) \\
    & \le (y-x_{k-1})\cdot \inf f([x_{k-1},y]) + (x_k - y)\cdot \inf f([y,x_k])
  \end{align*}
  and thus $S_*(f,P) \le S_*(f,P')$.
  Similarly, we obtain $S^*(f,P) \ge S^*(f,P')$.
  Therefore, \eqref{eq:39543} is proven.
  Now, if $P$ and $Q$ are any partitions, we observe that $P\cup Q$ 
  can be obtained from $P$ by adding one point at a time from $Q$. 
  Iterating \eqref{eq:39543}, we obtain \eqref{eq:5112309}.
  Taking the infimum over $Q$
  yields $S_*(f,P) \le I^*(f)$, and taking the supremum over 
  $P$ yields $I_*(f) \le I^*(f)$.
  \end{proof}


\begin{theorem}[integrability criteria]
\label{th:criteri_integrabilita}%
\mymark{*}%
\mymargin{integrability criteria}%
\index{criterion!of integrability}%
Let $f\colon[a,b]\to \RR$ be a bounded function.
\begin{enumerate}
\item
The function $f$ is Riemann-integrable if and only if
for every $\eps>0$ there exists a partition $P$
such that
\[
  S^*(f,P) - S_*(f,P) < \eps.
\]

\item
If $f$ is Riemann-integrable on $[a,b]$, then
there exists a sequence $P_n$ of partitions such that
\begin{equation}\label{eq:93765}
  \lim_{n\to +\infty} S^*(f,P_n)
  = \lim_{n\to+\infty} S_*(f,P_n)
  = \int_a^b f.
\end{equation}
Conversely, if $f$ is bounded and there exists a sequence $P_n$ 
of partitions of $[a,b]$ such that
\begin{equation*}
  \lim_{n\to +\infty} \enclose{S^*(f,P_n) - S_*(f,P_n)} = 0
\end{equation*}
then the function $f$ is Riemann-integrable and \eqref{eq:93765} holds.
\end{enumerate}
\end{theorem}
%
\begin{proof}
If there exists a partition $P$ such that $S^*(f,P)-S_*(f,P) < \eps$, we can
immediately conclude that
\[
I^*(f) - I_*(f) \le S^*(f,P) - S_*(f,P) < \eps.
\]
If this is true for every $\eps >0$, we deduce that $I^*(f) - I_*(f) = 0$ and
thus that $f$ is Riemann-integrable.

Conversely, for any $f$, by the properties
of $\sup$ and $\inf$,
there exist partitions $Q$ and $R$ such that
\[
  I^*(f) \ge S^*(f,Q) - \frac\eps 2
  \qquad\text{and}\qquad
  I_*(f) \le S_*(f,R) + \frac\eps 2
\]
from which, by the previous lemma, setting $P=Q\cup R$,
if $f$ is Riemann-integrable,
we obtain
\begin{align*}
S^*(f,P)-S_*(f,P) 
&\le S^*(f,Q) - S_*(f,R) \\
&\le I^*(f) + \frac\eps 2 - \enclose{I_*(f) - \frac \eps 2} = \eps.
\end{align*}

Now, let us prove the second point.
Suppose first that $f$ is Riemann-integrable on $[a,b]$.
Then, by the previous point, for every $n\in \NN$, setting $\eps=1/n$, we can
find a partition $P_n$ such that
\[
  S^*(f,P_n) - S_*(f,P_n) < \frac 1 n
\]
from which
\[
  I^*(f) \le S^*(f,P_n) \le S_*(f,P_n) + \frac 1 n
   \le I_*(f) + \frac 1 n
\]
therefore, taking the limit as $n\to +\infty$,
since $I^*(f) = I_*(f) = \int_a^b f$, it must hold that
\[
  \lim S^*(f,P_n) = \lim S_*(f,P_n) = \int_a^b f.
\]

Conversely, if
\[
 \lim_{n\to +\infty} S^*(f,P_n) - S_*(f,P_n) = 0
\]
for every $\eps>0$, there exists $n$ such that
\[
  S^*(f,P_n) - S_*(f,P_n) < \eps.
\]
By the previous point, we conclude that $f$ is Riemann-integrable.
On the other hand, we know that
\[
  S_*(f,P_n) \le I_*(f) = \int_a^b f = I^*(f) \le S^*(f,P_n)
\]
thus, if $S^*(f,P_n) - S_*(f,P_n) \to 0$, necessarily
the integral coincides with the limits of $S^*(f,P_n)$ and of
$S_*(f,P_n)$.
\end{proof}

The following exercise was carried out by Archimedes (287 BC - 212 BC) 
to calculate the area of a parabolic sector (method of exhaustion).

\begin{example}[calculation of the integral by partitions]%
\label{ex:integrale_quadrato}%
We show that for every $b>0$, the function $f(x)=x^2$ is Riemann-integrable on
the interval $[0,b]$ and we have
\[
 \int_0^b x^2\, dx = \frac{b^3}{3}.
\]
\end{example}
\begin{proof}
Consider the \emph{equally spaced} partitions of the interval $[0,b]$, that is,
divide $[0,b]$ into $N$ intervals each of width $b/N$:
\[
P_N = \ENCLOSE{\frac{kb}{N}\colon k \in 0, 1, \dots, N}.
\]
We have
\[
  S^*(f,P_N) 
      = \sum_{k=1}^N \frac b N \sup f\enclose{\Enclose{\frac{k-1}{N} b,\frac k N
   b}}
   = \frac{b}{N} \sum_{k=1}^N \frac{k^2b^2}{N^2}
   = \frac{b^3}{N^3} \sum_{k=1}^N k^2.
\]
Recall that (see exercise~\ref{ex:somma_quadrati}):
\[
  \sum_{k=1}^N k^2 = \frac{N(N+1)(2N+1)}{6} = \frac{2N^3+3N^2+N}{6}
  \sim \frac 1 3 N^3,\qquad \text{as $N\to +\infty$}.
\]
Thus, we have
\[
  S^*(f,P_N) \to \frac {b^3} 3, \qquad \text{as $N\to +\infty$}.
\]
Similarly, we find
\[
  S_*(f,P_N) 
    = \sum_{k=1}^N \frac{b}{N}\inf f\enclose{\Enclose{\frac{k-1} N b,\frac k N
  b}}
  = \frac{b}{N}\sum_{k=1}^N \frac{(k-1)^2b^2}{N^2}
  = \frac{b^3}{N^3} \sum_{k=0}^{N-1} k^2
\]
and thus 
\[
    S_*(f,P_N) \sim \frac {b^3} {N^3} \frac {(N-1)^3} 3 \to \frac {b^3} 3,
  \qquad \text{as $N\to +\infty$}.
\]
The proof is then concluded by applying
the criterion \eqref{eq:93765} of the previous
theorem~\ref{th:criteri_integrabilita}.
\end{proof}

\begin{theorem}[integral of a constant]
\label{th:integrale_costante}
If $f\colon[a,b]\to \RR$ is constant: $f(x) = c$, then
$f$ is Riemann-integrable and we have
\[
  \int_a^b f = c\cdot (b-a).
\]
\end{theorem}
%
\begin{proof}
Since on every $A\subset [a,b]$ we have
\[
  \sup_A f = \inf_A f = c
\]
it is easy to verify that
\[
  S^*(f,P) = S_*(f,P) = c\cdot (b-a)
\]
for any partition $P$ of $[a,b]$. The result follows immediately.
\end{proof}

Not all functions are Riemann-integrable, as shown by the following example.
\begin{example}[Dirichlet function]
\mymark{**}
\mymargin{Dirichlet function}
\index{function!of Dirichlet}
\index{integrability!counterexample}
Let $a<b$ and let $f\colon[a,b]\to \RR$ be the function defined by
\[
 f(x) =
 \begin{cases}
   1 & \text{if $x\in \QQ$}\\
   0 & \text{if $x\not \in \QQ$}.
 \end{cases}
\]
Then $f$ is not Riemann-integrable.
\end{example}
%
\begin{proof}
\mymark{*}
Let $P=\ENCLOSE{x_0, x_1, \dots, x_N}$ with $a=x_0 < x_1 < \dots < x_N = b$
be any partition of $[a,b]$.
Then it suffices to observe that, due to the density of rationals
(theorem~\ref{th:densita_frazioni}),
and of irrationals (exercise~\ref{ex:densita_irrazionali}),
in any subinterval $I=[x_{k-1}, x_k]$, there are infinitely many rational points
and infinitely many irrational points. Thus $\sup f(I)=1$ and $\inf f(I)=0$, and
consequently
\begin{align*}
  S^*(f,P) &= \sum_{k=1}^N (x_k - x_{k-1})\cdot 1 = b-a \\
  S_*(f,P) &= \sum_{k=1}^N (x_k - x_{k-1})\cdot 0 = 0
\end{align*}
from which $I^*(f) = b-a \neq 0 = I_*(f)$.
\end{proof}